\documentclass[12pt]{article}
\usepackage{amsmath,amssymb,amsthm}
\usepackage[margin=1in]{geometry}

\newtheorem{theorem}{Theorem}
\newtheorem{lemma}[theorem]{Lemma}
\newtheorem{corollary}[theorem]{Corollary}

\title{Problem 282: The Ackermann Function}
\author{Project Euler Solution}
\date{}

\begin{document}
\maketitle

\section{Problem Statement}

Define $A(n) = \operatorname{ack}(n,n)$ where $\operatorname{ack}$ is the Ackermann function:
\begin{align*}
\operatorname{ack}(0,m) &= m+1, \\
\operatorname{ack}(n,0) &= \operatorname{ack}(n-1,1) \quad (n>0), \\
\operatorname{ack}(n,m) &= \operatorname{ack}(n-1,\operatorname{ack}(n,m-1)) \quad (n,m>0).
\end{align*}
Find $\displaystyle\sum_{n=0}^{6} A(n) \pmod{14^8}$.

\section{Mathematical Development}

\begin{theorem}[Closed Forms for Levels 0--4]
For all $m \ge 0$:
\begin{enumerate}
\item $\operatorname{ack}(0,m) = m+1$,
\item $\operatorname{ack}(1,m) = m+2$,
\item $\operatorname{ack}(2,m) = 2m+3$,
\item $\operatorname{ack}(3,m) = 2^{m+3}-3$,
\item $\operatorname{ack}(4,m) = {}^{m+3}2 - 3$ (tower of $2$s of height $m+3$).
\end{enumerate}
\end{theorem}

\begin{proof}
Each level by induction on $m$.

(1)~Definition. (2)~Base: $\operatorname{ack}(1,0) = \operatorname{ack}(0,1) = 2$. Step: $\operatorname{ack}(1,m) = \operatorname{ack}(0,\operatorname{ack}(1,m-1)) = \operatorname{ack}(1,m-1)+1$, giving $m+2$ by induction.

(3)~Base: $\operatorname{ack}(2,0) = \operatorname{ack}(1,1) = 3$. Step: $\operatorname{ack}(2,m) = \operatorname{ack}(1,\operatorname{ack}(2,m-1)) = \operatorname{ack}(2,m-1)+2 = 2m+3$.

(4)~Base: $\operatorname{ack}(3,0) = \operatorname{ack}(2,1) = 5 = 2^3-3$. Step: $\operatorname{ack}(3,m) = 2\operatorname{ack}(3,m-1)+3 = 2(2^{m+2}-3)+3 = 2^{m+3}-3$.

(5)~Base: $\operatorname{ack}(4,0) = 2^4-3 = 13 = {}^{3}2 - 3$. Step: $\operatorname{ack}(4,m) = 2^{\operatorname{ack}(4,m-1)+3}-3 = 2^{{}^{m+2}2}-3 = {}^{m+3}2-3$.
\end{proof}

\begin{corollary}
$A(0)=1$, $A(1)=3$, $A(2)=7$, $A(3)=61$, $A(4) = {}^{7}2 - 3$.
\end{corollary}

\begin{theorem}[Modular Stabilization of Towers]
For any positive integer $M$, there exists $H \ge 0$ such that ${}^{h}2 \equiv {}^{H}2 \pmod{M}$ for all $h \ge H$. The minimal $H$ equals the length of the Carmichael chain $M, \lambda(M), \lambda(\lambda(M)), \ldots, 1$.
\end{theorem}

\begin{proof}
By strong induction on $M$. The case $M=1$ is trivial with $H=0$.

If $\gcd(2,M)=1$: by the generalized Euler theorem, $2^a \equiv 2^b \pmod{M}$ whenever $a \equiv b \pmod{\lambda(M)}$. Since $\lambda(M) < M$, the inductive hypothesis gives stabilization of the exponent ${}^{h-1}2 \bmod \lambda(M)$ for $h-1 \ge H'$, whence ${}^{h}2 \bmod M$ stabilizes for $h \ge H'+1$.

If $2 \mid M$: write $M = 2^a M'$ with $\gcd(2,M')=1$. For $h \ge a$, ${}^{h}2 \equiv 0 \pmod{2^a}$, and ${}^{h}2 \bmod M'$ stabilizes by the coprime case. By CRT, ${}^{h}2 \bmod M$ stabilizes.
\end{proof}

\begin{theorem}[CRT Decomposition for $M = 14^8$]
Let $M = 2^8 \cdot 7^8$. For $k \ge 4$, $A(k) \equiv 253 \pmod{256}$. The residue $A(k) \bmod 7^8$ is computed via the Carmichael chain, stabilizing for $k \ge 5$. CRT combines the residues.
\end{theorem}

\begin{proof}
Modulo $256$: for tower heights $h \ge 4$, ${}^{h}2 \equiv 0 \pmod{256}$, so $A(k) = {}^{(\cdots)}2 - 3 \equiv -3 \equiv 253$.

Modulo $7^8$: $\gcd(2,7^8)=1$, so we recurse down the Carmichael chain. For $A(5)$, the inner tower height exceeds $65536$, which vastly surpasses the chain length ($O(\log M) \approx 20$). Hence $A(5) \equiv A(6) \pmod{7^8}$: both equal the stabilized residue.
\end{proof}

\section{Algorithm}

\begin{verbatim}
small = [1, 3, 7, 61]
total = sum(small) mod M
a4 = (tower2_mod(7, M) - 3) mod M
stable = (tower2_mod(200, M) - 3) mod M
total += a4 + 2 * stable
return total mod M
\end{verbatim}

\section{Complexity Analysis}

\begin{itemize}
\item \textbf{Time:} $O(\log^2 M)$. The Carmichael chain has $O(\log M)$ levels, each requiring $O(\log M)$ work for modular exponentiation.
\item \textbf{Space:} $O(\log M)$ for the recursion stack.
\end{itemize}

\section{Answer}

\[
\boxed{1098988351}
\]

\end{document}
